\section{Przykłady przekrojowe: jedno zadanie, wiele kategorii}\label{sec:przekrojowe}

Zestawmy teraz w jednym miejscu kilka konstrukcji, żeby wyraźniej zobaczyć, co jest wspólne, a co zależy od kategorii.

\begin{przyklad}[Co jest iloczynem dwóch obiektów?]\label{cat:ex:cross-products}
\leavevmode\par\noindent
\begin{itemize}
\item W $\Set$: $A\times B$ i zwykłe rzuty.
\item W $\Grp$: $G\times H$ z mnożeniem współrzędnym.
\item W $\Top$: $X\times Y$ z topologią produktową.
\item W $\Vect_k$: $V\oplus W$ (dla dwóch czynników).
\item W posetach: $a\wedge b$ (kres dolny), jeśli istnieje.
\end{itemize}
We wszystkich przypadkach testem jest \emph{ta sama} własność uniwersalna z definicji~\ref{def:produkt}. Dla morfizmów $f:Z\to A$, $g:Z\to B$ do czynników w $\Set$ jedynym morfizmem o zadanych składowych jest $z\mapsto(f(z),g(z))$. W $\Grp$ jest on homomorfizmem, bo mnożenie jest współrzędne; w $\Vect_k$ jest liniowy z tego samego powodu dla dodawania i mnożenia przez skalar. W $\Top$ przeciwobraz bazowego zbioru otwartego $U\times V$ wynosi $f^{-1}(U)\cap g^{-1}(V)$, więc odwzorowanie jest ciągłe. Jednoznaczność we wszystkich czterech przypadkach wynika z jednoznaczności współrzędnych. W kategorii zbioru częściowo uporządkowanego warunek istnienia morfizmów $z\to a$ i $z\to b$ znaczy $z\leq a$ i $z\leq b$; istnienie jedynego morfizmu $z\to a\wedge b$ jest dokładnie własnością największego ograniczenia dolnego.
\end{przyklad}

\begin{przyklad}[Co jest koproduktem dwóch obiektów?]\label{cat:ex:cross-coproducts}
\leavevmode\par\noindent
\begin{itemize}
\item W $\Set$: suma rozłączna $A\sqcup B$.
\item W $\Grp$: iloczyn wolny $G*H$.
\item W $\Top$: suma rozłączna z topologią sumy topologicznej.
\item W $\Vect_k$: $V\oplus W$.
\item W posetach: $a\vee b$ (kres górny), jeśli istnieje.
\end{itemize}
Dla morfizmów $f:A\to Z$, $g:B\to Z$ w $\Set$ określamy jedyny morfizm na sumie rozłącznej, stosując $f$ na pierwszym składniku i $g$ na drugim. W $\Top$ przeciwobraz zbioru otwartego jest otwarty na obu składnikach, czyli otwarty w topologii sumy. W $\Vect_k$ jedynym liniowym morfizmem o zadanych ograniczeniach jest $(a,b)\mapsto f(a)+g(b)$, ponieważ $(a,b)=(a,0)+(0,b)$. Dla posetu istnienie morfizmów $a\to z$ i $b\to z$ jest równoważne $a\vee b\leq z$; jednoznaczność wynika z definicji kategorii posetu.

W $\Grp$ dla iloczynu wolnego $G*H$ przedłużenie homomorfizmów $f:G\to K$, $g:H\to K$ otrzymujemy, zastępując kolejne litery z $G$ lub $H$ ich obrazami i mnożąc w $K$. Redukcje słów (usuwanie jednostek i mnożenie sąsiednich liter z tego samego czynnika) zachowują wartość, więc ewaluacja jest homomorfizmem; jednoznaczność wynika z generowania przez oba czynniki. Dla porównania homomorfizm $G\times H\to K$ o tych samych ograniczeniach do czynników istnieje wtedy i tylko wtedy, gdy $f(x)g(y)=g(y)f(x)$ dla wszystkich $x\in G$, $y\in H$. Konieczność wynika z przemienności $(x,e)$ i $(e,y)$ w iloczynie; przy tym warunku wzór $(x,y)\mapsto f(x)g(y)$ zachowuje mnożenie i jest jedynym możliwym przedłużeniem. Zatem iloczyn i koprodukt nie mają tej samej własności w $\Grp$.
\end{przyklad}

\begin{przyklad}[Co oznacza strzałka?]
W $\Set$ strzałką jest funkcja; w $\Grp$ homomorfizm; w $\Rel$ relacja; w kategorii zbioru częściowo uporządkowanego jedyny morfizm $x\to y$ istnieje dokładnie wtedy, gdy $x\leq y$ (nie rozróżniamy dowodów nierówności); w kategorii monoidu z jednym obiektem strzałką jest element monoidu. Zatem nie ma sensu pytać, czy morfizm kategorii jest \emph{zawsze ciągły} lub \emph{zawsze różnowartościowy}: właściwe pytanie zależy od wybranej kategorii.
\end{przyklad}

\begin{przyklad}[Rozwiązywanie problemu przez własność uniwersalną]
Załóżmy, że chcemy określić homomorfizm $h:G*H\to K$. W podejściu elementarnym musielibyśmy opisać, jak $h$ działa na zredukowanych słowach, i sprawdzić zgodność z mnożeniem. W podejściu kategorialnym wystarczy podać dwa homomorfizmy $f:G\to K$ i $g:H\to K$: własność koproduktu gwarantuje, że istnieje dokładnie jeden $h$ z $h\circ i_G=f$, $h\circ i_H=g$. To jest praktyczna oszczędność pracy, a nie jedynie zmiana notacji.
\end{przyklad}

\begin{przyklad}[Funktor może ujawniać strukturę ukrytą]
Każda grupa $G$ działa na sobie przez lewe mnożenie. Traktując $G$ jako kategorię z jednym obiektem, ze składaniem $g\circ h=gh$, otrzymujemy funktor do $\Set$, który przypisuje jedynemu obiektowi zbiór $G$, a morfizmowi $g\in G$ permutację $L_g(x)=gx$. Mamy $L_e=\id_G$ i $L_g\circ L_h=L_{gh}$, a odwrotnością $L_g$ jest $L_{g^{-1}}$. Ta obserwacja prowadzi do twierdzenia Cayleya: odwzorowanie $g\mapsto L_g$ jest homomorfizmem $G\to\operatorname{Sym}(G)$, a jest różnowartościowe, ponieważ $L_g(e)=g$.
\end{przyklad}


\begin{przyklad}[Dlaczego własność uniwersalna daje jednoznaczność konstrukcji]
Załóżmy, że obiekt $P$ z rzutami $p_A,p_B$ ma uniwersalną własność iloczynu, podobnie jak $Q$ z rzutami $q_A,q_B$. Nie wybieramy żadnych reprezentantów elementów. Uniwersalność $Q$ daje jedyny $u:P\to Q$ z $q_Au=p_A$, $q_Bu=p_B$, a uniwersalność $P$ daje jedyny $v:Q\to P$ z $p_Av=q_A$, $p_Bv=q_B$. Wtedy $p_Avu=q_Au=p_A$, $p_Bvu=q_Bu=p_B$, więc $vu=\id_P$ z jednoznaczności morfizmu do $P$. Podobnie $q_Auv=p_Av=q_A$ i $q_Buv=p_Bv=q_B$, stąd $uv=\id_Q$. Izomorfizm $u$ jest jedyny wśród morfizmów zachowujących rzuty. Ten dowód jest typowym wzorcem: \emph{aby porównać dwa obiekty uniwersalne, zastosuj ich uniwersalności do siebie nawzajem}.
\end{przyklad}

\begin{przyklad}[Operacje na zbiorach a operacje na przestrzeniach liniowych]
Dla zbioru $X$ nie ma naturalnego działania $X\times X\to X$ wyznaczonego samą strukturą zbioru poza wyborem pierwszej albo drugiej współrzędnej. W istocie lemat Yonedy z przykładu~\ref{przyklad:yoneda-tuple} daje dokładnie dwa naturalne odwzorowania $X^2\to X$: $(x,y)\mapsto x$ i $(x,y)\mapsto y$. Natomiast dla przestrzeni wektorowej $V$ istnieje naturalne dodawanie $V\oplus V\to V$, $(v,w)\mapsto v+w$, ponieważ każdy liniowy $f:V\to W$ spełnia $f(v+w)=f(v)+f(w)$. \emph{Naturalność zawsze zależy od kategorii i jej morfizmów}, a nie tylko od wzoru na elementach.
\end{przyklad}


\begin{przyklad}[Iloczyn i koprodukt w kategorii relacji]\label{cat:ex:rel-biproduct}
W $\Rel$ morfizm $R:A\to B$ jest podzbiorem $A\times B$, a składanie jest składaniem relacji. Suma rozłączna $P=A\sqcup B$ jest tu zarówno iloczynem, jak i koproduktem. Niech $i_A,i_B$ będą relacjami będącymi grafami włączeń, a $p_A,p_B$ relacjami odwrotnymi do tych grafów (nie funkcjami na całym $P$).

Dla relacji $R:Z\to A$, $S:Z\to B$ definiujemy $H:Z\to P$ przez
\[
z\,H\,i_A(a)\ \Longleftrightarrow\ z\,R\,a,\qquad
z\,H\,i_B(b)\ \Longleftrightarrow\ z\,S\,b,
\]
gdzie $i_A(a),i_B(b)$ oznaczają oznakowane elementy sumy. Z definicji składania $p_A\circ H=R$, $p_B\circ H=S$. Warunki te wyznaczają każdą parę należącą do $H$, ponieważ każdy element $P$ należy do dokładnie jednego składnika. To dowodzi własności iloczynu. Odwrotnie, dla $R':A\to Z$, $S':B\to Z$ określamy $H':P\to Z$ warunkami $i_A(a)\,H'\,z$ wtedy i tylko wtedy, gdy $a\,R'\,z$, oraz $i_B(b)\,H'\,z$ wtedy i tylko wtedy, gdy $b\,S'\,z$. Otrzymujemy jedyną relację spełniającą $H'i_A=R'$, $H'i_B=S'$, czyli własność koproduktu. Argument obejmuje również puste składniki. W szczególności iloczyn dwóch singletonów w $\Rel$ ma dwa elementy, a nie jeden jak w $\Set$.
\end{przyklad}

\begin{przyklad}[Klasyfikacja naturalnych operacji liniowych]\label{cat:ex:linear-natural-operations}
Ustalmy ciało $k$. Każda transformacja naturalna $\tau:V\oplus V\Rightarrow V$ (od funktora $V\mapsto V\oplus V$ do identyczności na $\Vect_k$) ma postać
\[
\tau_V(v,w)=av+bw
\]
dla jednoznacznie wyznaczonych skalarów $a,b\in k$. Istotnie, niech $a=\tau_k(1,0)$ i $b=\tau_k(0,1)$. Dla dowolnego $v\in V$ liniowy morfizm $f_v:k\to V$, $t\mapsto tv$, daje z naturalności
$\tau_V(v,0)=f_v(\tau_k(1,0))=av$. Analogicznie $\tau_V(0,w)=bw$, a liniowość $\tau_V$ daje żądany wzór. Odwrotnie, każda para skalarów definiuje liniowe składowe, a dla każdego liniowego $f:V\to W$ zachodzi $f(av+bw)=af(v)+bf(w)$. Jednoznaczność współczynników wynika z wartości na $(1,0)$ i $(0,1)$ w $k\oplus k$. Zatem w $\Vect_k$ naturalne operacje dwuargumentowe tworzą rodzinę parametryzowaną przez $k^2$, w przeciwieństwie do dwóch projekcji w $\Set$.
\end{przyklad}
